


\newgeometry{margin=0pt}
\thispagestyle{empty}

\begin{tikzpicture}[remember picture, overlay]

  % Baggrund
  \fill[jagtgroen]
    (current page.south west) rectangle (current page.north east);

  % Ydre ramme
  \draw[jagtguld, line width=2.5pt]
    ([shift={(1.2cm,1.2cm)}]current page.south west) rectangle
    ([shift={(-1.2cm,-1.2cm)}]current page.north east);

  % Indre ramme
  \draw[jagtguld, line width=0.8pt]
    ([shift={(1.6cm,1.6cm)}]current page.south west) rectangle
    ([shift={(-1.6cm,-1.6cm)}]current page.north east);

  % Titel
  \node[text=white, align=center, text width=0.72\paperwidth, font=\fontsize{52}{60}\selectfont\bfseries]
    at ([yshift=3.8cm]current page.center) {Jagtprøven};

  % Streg under titel
  \draw[jagtguld, line width=1.5pt]
    ([shift={(-4cm,1.1cm)}]current page.center) --
    ([shift={(4cm,1.1cm)}]current page.center);

  % Undertitel
  \node[text=white, align=center, text width=0.68\paperwidth, font=\fontsize{22}{28}\selectfont]
    at ([yshift=-0.6cm]current page.center) {Vejledning til effektiv beståelse af den teoretiske og praktiske del af jagtprøven};

  % Beskrivelse
  \node[text=jagtguld, align=center, text width=0.7\paperwidth, font=\fontsize{14.5}{19}\selectfont\itshape]
    at ([yshift=-3.6cm]current page.center) {Hent seneste version gratis via osjagttegn.dk/bog};

  % Open Source
  \node[text=jagtguld, align=center, text width=0.7\paperwidth, font=\fontsize{13.5}{17.5}\selectfont\itshape]
    at ([yshift=-5.1cm]current page.center) {Lærebogen er Danmarks første Open Source introduktionsbog til jagttegn};

  % Licens
  \node[text=jagtguld, align=center, text width=0.7\paperwidth, font=\fontsize{13.5}{17.5}\selectfont\itshape]
    at ([yshift=-6.5cm]current page.center) {Licenseret under CC BY 4.0};

  % Fod
  \node[text=white, align=center, font=\small]
    at ([yshift=2.6cm]current page.south) 
    {Aleksander Sønder | \href{https://osjagttegn.dk}};

\end{tikzpicture}

\clearpage
\restoregeometry